\chapter{Conclusioni}
\label{cap:Conclusioni}

\section{Sintesi dei contributi}

Questa tesi ha affrontato il problema dell'efficienza energetica nel
paradigma Function-as-a-Service a partire da una lacuna concreta: le
piattaforme serverless non dispongono, in generale, di meccanismi nativi
per adattare l'esecuzione delle funzioni al costo carbonico dell'energia
disponibile nel momento e nel luogo dell'invocazione. Nel comportamento
tradizionale, una funzione viene eseguita nella propria versione di
riferimento indipendentemente dal contesto energetico, senza distinguere
tra un periodo in cui la rete elettrica è alimentata prevalentemente da
fonti rinnovabili e uno in cui è dominata da fonti fossili.

Il contributo complessivo del lavoro consiste nell'estensione della
piattaforma Serverledge con un modello di esecuzione carbon-aware, nel
quale la scelta della variante non è fissata staticamente ma viene
determinata a runtime in funzione della carbon intensity corrente della zona di
esecuzione, dei profili energetici osservati e del livello di qualità
associato alle alternative disponibili. Questo contributo si articola in
quattro direzioni principali.

Il primo contributo è il modello di funzione logica con varianti. Una
funzione logica non corrisponde a una sola implementazione, ma a un insieme
di varianti compatibili dal punto di vista dell'interfaccia e dello scopo,
ma diverse per costo energetico e qualità del risultato. In questo modo il
trade-off energia--qualità diventa un elemento esplicito del sistema, e non
una proprietà implicita lasciata interamente allo sviluppatore.

Il secondo contributo riguarda il monitoraggio energetico a granularità di
funzione. Tramite l'integrazione di Kepler con Prometheus, il sistema
raccoglie metriche di consumo a livello di container e le associa alle
singole invocazioni tramite un worker periodico. I profili energetici
ottenuti vengono aggiornati nel tempo con una media mobile esponenziale e
memorizzati in etcd, rendendoli disponibili allo scheduler durante la
selezione. Questo meccanismo permette di superare una visione puramente
statica del consumo, adattando le stime al comportamento osservato sul nodo
corrente.

Il terzo contributo è l'integrazione dell'incertezza nella selezione. Nelle
fasi iniziali, quando alcune varianti sono state osservate poche volte, il
sistema applica una stima ottimistica ispirata al principio Lower Confidence
Bound. Le varianti meno campionate vengono così esplorate con maggiore
probabilità, riducendo il rischio che una stima iniziale imprecisa escluda
prematuramente alternative potenzialmente efficienti.

Il quarto contributo è lo scheduler carbon-aware. A partire dalla carbon
intensity della zona di esecuzione, il sistema calcola un parametro
$\lambda$ tramite una funzione sigmoide calibrata sulla distribuzione
storica locale. Tale parametro regola una scalarizzazione lineare della
frontiera di Pareto energia--qualità: valori elevati di $\lambda$
privilegiano varianti più accurate, mentre valori ridotti orientano la
selezione verso varianti più efficienti. La calibrazione locale consente di
interpretare la carbon intensity non come valore assoluto, ma come
posizione rispetto al comportamento tipico della rete elettrica considerata.

\section{Risultati principali}

La valutazione sperimentale ha considerato sei funzioni, tre numeriche e
tre di classificazione, su cinque zone europee con profili di carbon
intensity differenti. Gli esperimenti hanno confrontato lo scheduler
carbon-aware con due baseline statiche: $\Lambda_1$, che seleziona sempre
la variante a massima qualità, e $\Lambda_0$, che seleziona sempre la
variante a minima energia.

L'Esperimento~1 ha verificato il comportamento della logica decisionale con
una calibrazione europea comune. I risultati mostrano che lo scheduler si
sposta lungo la frontiera di Pareto in modo coerente con il segnale
carbonico: nelle zone a bassa carbon intensity tende a selezionare varianti
più accurate, mentre nelle zone più carbon-intensive privilegia varianti
più efficienti.

L'Esperimento~2 ha introdotto una calibrazione locale della sigmoide,
interpretando ogni valore di carbon intensity rispetto alla storia della
zona. Questa scelta produce un comportamento più equilibrato: anche zone con
CI media elevata possono generare decisioni intermedie quando il valore
corrente è ordinario rispetto alla propria distribuzione storica. Rispetto
alla baseline $\Lambda_1$, lo scheduler ottiene risparmi energetici compresi
tra il 62\% e il 92\% a seconda della funzione e della zona, con una
riduzione stimata delle emissioni di CO\textsubscript{2} fino al 94{,}82\%.
Il costo qualitativo rimane contenuto per le funzioni numeriche, con perdite
comprese tra circa 0{,}01\% e 0{,}61\%, mentre risulta più evidente nelle
funzioni ML, dove i passaggi tra modelli distinti producono riduzioni di
qualità comprese tra il 4{,}34\% e il 10{,}55\%. Rispetto alla baseline
$\Lambda_0$, lo scheduler mantiene invece un guadagno qualitativo compreso
tra 0{,}13\% e 4{,}62\% per le funzioni numeriche e tra 23{,}40\% e
32{,}19\% per le funzioni ML.

L'Esperimento~3 ha esteso la valutazione a un flusso di richieste
concorrenti con carico crescente. I risultati confermano che la presenza di
concorrenza non altera la semantica della decisione carbon-aware: lo
scheduler continua a posizionarsi tra le due baseline, riducendo energia ed
emissioni rispetto a $\Lambda_1$ e preservando qualità rispetto a
$\Lambda_0$. Inoltre, l'esperimento evidenzia un effetto operativo
ulteriore. La baseline $\Lambda_1$ raggiunge rapidamente il proprio limite
di capacità e si stabilizza attorno a 1{,}25~req/s, mentre
\emph{adaptive-ci} mantiene un throughput effettivo aderente al profilo
configurato, con una divergenza visibile solo alla pressione massima. La
scelta dinamica della variante contribuisce quindi non solo alla riduzione
dell'impatto carbonico, ma anche alla stabilità del sistema sotto carico
variabile.

Nel complesso, i risultati indicano che la riduzione delle emissioni non è
ottenuta degradando indistintamente la qualità del servizio, ma scegliendo
di volta in volta la variante più adatta al costo carbonico dell'esecuzione.
Il beneficio emerge quando il sistema dispone di alternative sufficientemente
differenziate e di un segnale ambientale capace di guidarne la scelta.

\section{Limiti del lavoro attuale}

Nonostante i risultati ottenuti, il sistema proposto presenta alcune
limitazioni che definiscono i confini di validità del contributo e indicano
le principali direzioni di miglioramento.

\paragraph{Varianti definite manualmente.}
Il sistema richiede che le varianti di ciascuna funzione siano progettate e
registrate a priori dallo sviluppatore. Non esiste, nel prototipo attuale,
un meccanismo automatico di generazione o validazione delle varianti. La
copertura del trade-off energia--qualità dipende quindi dalla qualità,
dalla varietà e dalla granularità delle implementazioni fornite. Questo
vincolo limita la scalabilità dell'approccio in contesti con molte funzioni
o con applicazioni aggiornate frequentemente.

\paragraph{Qualità modellata come attributo statico.}
La qualità associata alle varianti è trattata come un valore noto e
registrato nel sistema. Questa scelta è adeguata per una valutazione
controllata, ma non cattura possibili variazioni legate al tipo di input,
alla distribuzione dei dati o alla sensibilità applicativa della singola
richiesta. Una stessa perdita media di qualità può essere trascurabile per
alcune invocazioni e non accettabile per altre; il modello attuale non
dispone di un meccanismo per distinguere automaticamente questi casi.

\paragraph{Granularità discreta della frontiera di Pareto.}
La qualità del compromesso ottenibile dipende direttamente dalla densità
della frontiera di Pareto. Quando le varianti sono numerose e vicine, come
nelle funzioni numeriche parametrizzabili, lo scheduler può muoversi in modo
graduale tra qualità ed energia. Quando invece la frontiera è composta da
pochi modelli distinti, come in molte funzioni ML, il passaggio da una
variante all'altra produce salti qualitativi più marcati. Anche una logica
di selezione ottimale non può individuare compromessi intermedi se tali
punti non sono presenti tra le varianti disponibili.

\paragraph{Ambiente sperimentale controllato.}
La valutazione è stata condotta su un insieme limitato di funzioni
rappresentative e su un'infrastruttura di laboratorio basata su un singolo
nodo Serverledge. Il workload dell'Esperimento~3 introduce concorrenza e
carico variabile, ma resta deterministico e sintetico. I risultati
dimostrano la fattibilità del meccanismo, ma non esauriscono la varietà di
condizioni presenti in ambienti produttivi, dove intervengono autoscaling,
multi-tenancy, interferenze tra applicazioni, burst non prevedibili e
vincoli di latenza più stringenti.

\paragraph{Stima energetica e confini del calcolo carbonico.}
In ambienti privi di accesso diretto all'interfaccia RAPL, Kepler opera
tramite modelli di stima del consumo. Il rumore introdotto da tali modelli
può influenzare l'aggiornamento dei profili energetici, soprattutto per
funzioni a basso consumo assoluto. Questo limite non riguarda la natura
statica dei profili, che nel sistema vengono aggiornati online ed esplorati
nelle fasi iniziali, ma l'accuratezza della sorgente di misura che alimenta
tale aggiornamento. Inoltre, la stima delle emissioni considera solo la
quota operativa dell'esecuzione, cioè l'energia consumata dalla funzione
moltiplicata per la carbon intensity della zona. Non viene invece modellata
la \emph{embodied carbon}, ovvero la quota di emissioni associata alla
produzione, al trasporto, all'installazione e alla dismissione
dell'hardware, ammortizzata sul ciclo di vita dell'infrastruttura
computazionale~\cite{gupta2021chasingcarbon}. Questa scelta delimita il
perimetro del lavoro: lo scheduler riduce l'impatto carbonico operativo
delle invocazioni, ma non pretende di stimare l'intera impronta ambientale
del sistema.

\paragraph{Assenza di decisioni distribuite.}
Il sistema decide quale variante eseguire sul nodo corrente, ma non sceglie
dove eseguirla. In uno scenario cloud-edge reale, la selezione della
variante e il posizionamento geografico dell'esecuzione sarebbero decisioni
interdipendenti. Sarebbe necessario valutare se spostare l'esecuzione verso
una zona con carbon intensity più bassa consenta di ottenere un compromesso
migliore tra energia, qualità ed emissioni, tenendo però conto del costo di
trasferimento introdotto dalla decisione: latenza di rete, energia consumata
per inviare dati o stato applicativo e possibile overhead di avvio sul nodo
remoto. Questa interazione resta fuori dal perimetro del prototipo attuale.

\section{Sviluppi futuri}

I limiti individuati aprono diverse direzioni di ricerca e sviluppo, sia a
breve termine sia verso scenari cloud-edge più ampi.

\paragraph{Generazione automatica di varianti.}
Un'evoluzione naturale consiste nell'integrare strumenti per generare o
suggerire automaticamente varianti alternative. Tecniche di traduzione e
rifattorizzazione automatica del codice, come quelle esplorate in
\cite{refaas}, potrebbero produrre implementazioni più efficienti dal punto
di vista energetico. Nel caso delle funzioni ML, tecniche di compressione
del modello potrebbero svolgere un ruolo analogo: il \emph{pruning} rimuove
connessioni o componenti poco rilevanti, la quantizzazione rappresenta pesi
e attivazioni con precisione numerica ridotta, mentre la distillazione
addestra un modello più piccolo a imitare il comportamento di un modello più
grande~\cite{han2016deepcompression,hinton2015distilling}. Queste tecniche
permetterebbero di costruire famiglie di varianti con livelli progressivi
di costo e accuratezza. Una frontiera più densa ridurrebbe i salti tra
varianti e permetterebbe allo scheduler di modulare il compromesso con
maggiore continuità.

\paragraph{Validazione dinamica della qualità.}
Un'estensione importante riguarda il passaggio da valori statici di qualità
a metriche dipendenti dagli input e dal contesto applicativo. Lo scheduler
potrebbe integrare vincoli di qualità minima, classi di richieste con
criticità diversa o meccanismi di monitoraggio online della correttezza.
Questo permetterebbe di distinguere richieste per cui l'approssimazione è
accettabile da richieste per cui la piattaforma deve preservare la variante
più accurata.

\paragraph{Calibrazione adattiva e predittiva della carbon intensity.}
La sigmoide è attualmente calibrata su dati storici. Una calibrazione
periodica basata su finestre mobili permetterebbe di adattarsi a variazioni
strutturali del mix energetico locale. Un passo ulteriore consiste
nell'integrare previsioni di carbon intensity, così da pianificare le scelte
non solo rispetto al valore corrente ma anche rispetto all'evoluzione
attesa della rete elettrica.

\paragraph{Scheduling distribuito multi-nodo.}
In un deployment cloud-edge federato, lo scheduler dovrebbe decidere sia la
variante sia il nodo di esecuzione. Questo richiede di combinare carbon
intensity, latenza di rete, capacità residua, stato dei container, costo
energetico della variante e costo di trasferimento. L'obiettivo non sarebbe
spostare sempre l'esecuzione verso la zona più pulita, ma verificare quando
il beneficio di una minore carbon intensity compensa l'energia e la latenza
aggiuntive introdotte dal trasferimento. L'estensione multi-nodo
trasformerebbe quindi il modello da selezione locale della variante a
problema congiunto di routing e scheduling carbon-aware.

\paragraph{Integrazione con politiche di SLA e pricing.}
La selezione carbon-aware introduce una differenziazione del servizio: a
parità di funzione logica, l'utente può ricevere varianti diverse in base
al contesto ambientale e alle preferenze espresse. Uno sviluppo futuro
consiste nel rendere questa scelta parte del contratto di servizio,
permettendo all'utente di specificare vincoli su qualità minima, latenza,
impatto carbonico massimo o costo economico. In un modello pay-as-you-go,
tali preferenze potrebbero essere riflesse anche in politiche di pricing
differenziate.

\paragraph{Valutazione su workload e infrastrutture reali.}
Infine, il sistema dovrebbe essere validato su deployment multi-nodo e su
tracce di workload reali, includendo scenari di autoscaling, cold start,
interferenza tra tenant e variazioni rapide della domanda. Una campagna
sperimentale di questo tipo permetterebbe di misurare il comportamento
dello scheduler in condizioni più vicine alla produzione e di quantificare
il costo operativo dei meccanismi introdotti.
